% Generated draft: requires mathematical and source review.
\documentclass[11pt]{article}
\usepackage[margin=25mm]{geometry}
\usepackage{fontspec}
\setmainfont{Latin Modern Roman}
\usepackage{amsmath,amssymb,mathtools,unicode-math}
\setmathfont{Latin Modern Math}
\usepackage{enumitem,needspace,xurl,hyperref}
\setlength{\parindent}{0pt}
\setlength{\parskip}{5pt}
\setlength{\emergencystretch}{3em}
\newcommand{\N}{\mathbb{N}}
\newcommand{\Z}{\mathbb{Z}}
\newcommand{\Q}{\mathbb{Q}}
\newcommand{\R}{\mathbb{R}}
\newcommand{\C}{\mathbb{C}}
\newcommand{\F}{\mathbb{F}}
\newcommand{\A}{\mathbb{A}}
\newcommand{\PP}{\mathbb{P}}
\newcommand{\E}{\mathbb{E}}
\newcommand{\eps}{\varepsilon}
\newcommand{\dd}{\,\mathrm{d}}
\begin{document}
\section*{UnsolvedMath ID: 30000122}
\textbf{Source problem code:} catalogue entry
\subsection{Definitions and mathematical statement}
\textbf{Classification question.} For which integers \(n\) does the odd-dimensional sphere \(S^{2n-1}\) admit a foliation by complex manifolds? Here \(S^{2n-1}=\{x\in\R^{2n}:\|x\|=1\}\). The source does not state the leaf dimension, codimension, regularity, or whether singular foliations are permitted. Allowing zero-dimensional leaves gives the foliation by points on every sphere, so a nontrivial leaf-dimension requirement is essential.\par Under the usual interpretation suggested by the cited Reeb-foliation context, take \(n\ge2\) and ask for a regular real-codimension-one foliation \(\mathcal F\) on \(S^{2n-1}\), together with a smoothly varying complex structure \(J\) on its rank-\(2n-2\) tangent bundle \(T\mathcal F\), such that \(J^2=-\operatorname{id}\) and each leaf with the induced structure is a complex manifold of complex dimension \(n-1\). Thus leafwise integrability is required; a leafwise almost complex structure alone is insufficient. Equivalently, one seeks a smooth foliation whose leafwise coordinate changes are holomorphic in the complex leaf variables, with the intended transverse regularity specified. This is a conditional codimension-one formulation, not a claim that the missing source explicitly imposes it. If other positive leaf dimensions or regularities are intended, they define different existence questions. No complex-manifold structure on the entire odd-dimensional sphere is being requested.\par The supplied literature triage checked 2026-08-21 labels the classification open. It cites low-dimensional examples such as the Reeb foliation of \(S^3\), and reports that the previously claimed \(S^5\) construction was corrected and does not establish such a foliation of \(S^5\). Failure of that construction is not an impossibility theorem for \(S^5\). These dated claims and the exact foliation hypotheses of the cited correction require independent verification; no current classification is asserted here.
\subsection{Short English statement}
Determine which odd-dimensional spheres carry a nontrivial foliation with complex leaves, under the intended dimension and regularity conventions. The likely codimension-one interpretation is explicit. The supplied correction invalidates a claimed \(S^5\) construction; it does not prove nonexistence.
\subsection{Sources}
\begin{itemize}\item Local UnsolvedMath registry, permanent ID 30000122, code catalogue entry; \emph{Complex Foliations of Odd-Dimensional Spheres}, Oberwolfach Reports Open Problems.\item \emph{Komplexe Analysis (2005)}, Oberwolfach Reports, DOI 10.4171/owr/2004/42, pp. 2171--2216. \url{https://doi.org/10.4171/owr/2004/42}. Proposed year and year tag: 2004; no proposer named.\item Supplied literature label \emph{Meersseman--Verjovsky correction}: \url{https://annals.math.princeton.edu/2011/174-3/p11}.\item Supplied literature label \emph{OWR report}: \url{https://doi.org/10.4171/owr/2004/42}.\item Literature triage checked 2026-08-21, status \texttt{open}; record updated 2026-09-16. The \texttt{corrected\_verified} field concerns balanced notation and removal of a repeated question label. Verify the dated literature claims independently.\end{itemize}
\end{document}
